\documentclass[a4paper,12pt,twoside,parskip=half]{scrartcl}
\usepackage[utf8]{inputenc}
\usepackage[
    module={Tabellen},
    link={LaTeX.dlrg.de}
]{dlrg}
\begin{document}

\section{Allgemeiner Bereich}

Dieses ist ein Beispiel für ein einfaches Dokument mit Bauchbinde. Es hat zusätzlich zur Demostation das Modul Tabellen eingebunden.

\section{Tabellen}
In diesem Bereich werden die verschiedenen Möglichkeiten für Tabellen dargestellt, die mit dem Modul Tabellen möglich sind.

Die Umgebung \textbf{dlrgTblr} liefert das einfache Ergebnis einer Tabelle, bei der alle Zellen schwarz umrandet sind.

\begin{dlrgTblr}{colspec={XXX}}
    Feld1 & Feld2 & Feld3 \\
    Inhalt 1 & Inhalt 2 & Ein viel längerer Inhalt als die anderen Elemente \\
    Inhalt 4 & Inhalt 5 & Inhalt 6\\
    Inhalt 7 & Inhalt 8 & \\
\end{dlrgTblr}

Durch die Option \verb|noHlines| werden die horizontalen Linien nicht erstellt. Analog funktioniert dieses mit \verb|noVlines| für die vertikalen Linien.

\begin{dlrgTblr}{colspec={XXX}}[noHlines]
    Feld1 & Feld2 & Feld3 \\
    Inhalt 1 & Inhalt 2 & Ein viel längerer Inhalt als die anderen Elemente \\
    Inhalt 4 & Inhalt 5 & Inhalt 6\\
    Inhalt 7 & Inhalt 8 & \\
\end{dlrgTblr}

Mit der Option \verb|ueberschrift| kann man die erste Zeile als Überschrift hervorheben. Dabei wird diese bei der einfachen Ansicht grau hinterlegt.

\begin{dlrgTblr}{colspec={XXX}}[ueberschrift]
    Feld1 & Feld2 & Feld3 \\
    Inhalt 1 & Inhalt 2 & Ein viel längerer Inhalt als die anderen Elemente \\
    Inhalt 4 & Inhalt 5 & Inhalt 6\\
    Inhalt 7 & Inhalt 8 & \\
\end{dlrgTblr}

Mit der Option \verb|layout=rot| werden die Linien der Tabelle in Rot dargestellt.

\begin{dlrgTblr}{colspec={XXX}}[layout=rot]
    Feld1 & Feld2 & Feld3 \\
    Inhalt 1 & Inhalt 2 & Ein viel längerer Inhalt als die anderen Elemente \\
    Inhalt 4 & Inhalt 5 & Inhalt 6\\
    Inhalt 7 & Inhalt 8 & \\
\end{dlrgTblr}

Bei der Kombination von \verb|ueberschrift| und \verb|layout=rot| kann man deutlich erkennen, dass die Darstellung der Überschrift auch an das rote Layout angepasst ist.

\begin{dlrgTblr}{colspec={XXX}}[ueberschrift,layout=rot]
    Feld1 & Feld2 & Feld3 \\
    Inhalt 1 & Inhalt 2 & Ein viel längerer Inhalt als die anderen Elemente \\
    Inhalt 4 & Inhalt 5 & Inhalt 6\\
    Inhalt 7 & Inhalt 8 & \\
\end{dlrgTblr}

\end{document}
